\documentclass[11pt,a4paper]{article}

% ---------------------------------------------------------------
%  Curriculum Vitae -- Stanisław Szufa
%  EPFL IC Faculty Search 2026
% ---------------------------------------------------------------
\usepackage[T1]{fontenc}
\usepackage[utf8]{inputenc}
\usepackage{mathpazo}
\renewcommand{\ttdefault}{lmtt}
\usepackage[hmargin=2.5cm,vmargin=2.2cm]{geometry}
\usepackage{microtype}
\usepackage{enumitem}
\usepackage{titlesec}
\usepackage{xcolor}
\usepackage{tabularx}
\usepackage{etoolbox}
\usepackage[hidelinks]{hyperref}
\hypersetup{pdfauthor={Stanisław Szufa},pdftitle={Stanisław Szufa: Curriculum Vitae}}

\definecolor{accent}{RGB}{160,20,30}   % same accent as the other documents
\definecolor{muted}{RGB}{90,90,90}
\titleformat{\section}{\large\bfseries\color{accent}}{}{0pt}{}[\vspace{1pt}{\color{accent!35}\titlerule[0.6pt]}]
\titlespacing*{\section}{0pt}{12pt}{5pt}
% Keep a section heading together with its first entries: start a new page
% if less than the given space is left.
\newcommand{\needspace}[1]{\par
  \ifdim\pagegoal<\maxdimen\ifdim\dimexpr\pagegoal-\pagetotal\relax<#1\relax\newpage\fi\fi}
\pretocmd{\section}{\needspace{6\baselineskip}}{}{}
\setlength{\parindent}{0pt}
\setlength{\parskip}{3pt}
\setlist{nosep,leftmargin=1.1em}
\pagestyle{plain}

% One dated entry: \entry{dates}{title}{place}{optional details}
\newlength{\datecol}\setlength{\datecol}{2.6cm}
\newcommand{\entry}[4]{%
  \par\noindent
  \begin{minipage}[t]{\datecol}\leavevmode{\color{muted}#1}\end{minipage}%
  \begin{minipage}[t]{\dimexpr\linewidth-\datecol\relax}\raggedright
    \textbf{#2}\ifblank{#3}{}{, #3}%
    \ifblank{#4}{}{\\[1pt]{\small\color{black!85}#4\par}}
  \end{minipage}\par\vspace{3pt}}
% One selected paper, with its label in the date column: \paper{label}{text}
\newcommand{\paper}[2]{%
  \par\noindent
  \begin{minipage}[t]{\datecol}\leavevmode{\color{muted}[#1]}\end{minipage}%
  \begin{minipage}[t]{\dimexpr\linewidth-\datecol\relax}\raggedright
    #2
  \end{minipage}\par\vspace{3pt}}

% ---------------------------------------------------------------
%  Shared document header -- Stanisław Szufa, EPFL application
%  Load after xcolor and hyperref, and after defining the accent colour:
%    % ---------------------------------------------------------------
%  Shared document header -- Stanisław Szufa, EPFL application
%  Load after xcolor and hyperref, and after defining the accent colour:
%    % ---------------------------------------------------------------
%  Shared document header -- Stanisław Szufa, EPFL application
%  Load after xcolor and hyperref, and after defining the accent colour:
%    \input{../shared/header}
%  Usage:
%    \docheader{Research Statement}        % name, title, contact line, rule
%    \docheader{}                          % no title (cover letter)
%    \docheader[extra contact line]{Title} % second contact line (CV)
%    \headercontactfalse before \docheader    % no contact lines
% ---------------------------------------------------------------
\providecolor{muted}{RGB}{90,90,90}
\newcommand{\headersep}{\enspace{\color{accent!60}\textbullet}\enspace}
% Contact lines are shown unless a document says \headercontactfalse.
\newif\ifheadercontact\headercontacttrue
\newcommand{\docheader}[2][]{%
  \begin{center}
    {\LARGE\bfseries\color{accent} Stanisław Szufa}\par
    \if\relax\detokenize{#2}\relax\else
      \vspace{3pt}{\large\color{muted} #2}\par
    \fi
    \ifheadercontact
    \vspace{5pt}
    {\small\color{muted}
    University of Geneva, Switzerland\par\vspace{1pt}
    \href{mailto:stanislaw.szufa@unige.ch}{stanislaw.szufa@unige.ch}\headersep
    +41 76 472 16 09%
    \if\relax\detokenize{#1}\relax\else\headersep#1\fi\par}
    \fi
  \end{center}
  \vspace{-6pt}
  {\color{accent!40}\rule{\linewidth}{0.5pt}}\par
  \vspace{2pt}}

%  Usage:
%    \docheader{Research Statement}        % name, title, contact line, rule
%    \docheader{}                          % no title (cover letter)
%    \docheader[extra contact line]{Title} % second contact line (CV)
%    \headercontactfalse before \docheader    % no contact lines
% ---------------------------------------------------------------
\providecolor{muted}{RGB}{90,90,90}
\newcommand{\headersep}{\enspace{\color{accent!60}\textbullet}\enspace}
% Contact lines are shown unless a document says \headercontactfalse.
\newif\ifheadercontact\headercontacttrue
\newcommand{\docheader}[2][]{%
  \begin{center}
    {\LARGE\bfseries\color{accent} Stanisław Szufa}\par
    \if\relax\detokenize{#2}\relax\else
      \vspace{3pt}{\large\color{muted} #2}\par
    \fi
    \ifheadercontact
    \vspace{5pt}
    {\small\color{muted}
    University of Geneva, Switzerland\par\vspace{1pt}
    \href{mailto:stanislaw.szufa@unige.ch}{stanislaw.szufa@unige.ch}\headersep
    +41 76 472 16 09%
    \if\relax\detokenize{#1}\relax\else\headersep#1\fi\par}
    \fi
  \end{center}
  \vspace{-6pt}
  {\color{accent!40}\rule{\linewidth}{0.5pt}}\par
  \vspace{2pt}}

%  Usage:
%    \docheader{Research Statement}        % name, title, contact line, rule
%    \docheader{}                          % no title (cover letter)
%    \docheader[extra contact line]{Title} % second contact line (CV)
%    \headercontactfalse before \docheader    % no contact lines
% ---------------------------------------------------------------
\providecolor{muted}{RGB}{90,90,90}
\newcommand{\headersep}{\enspace{\color{accent!60}\textbullet}\enspace}
% Contact lines are shown unless a document says \headercontactfalse.
\newif\ifheadercontact\headercontacttrue
\newcommand{\docheader}[2][]{%
  \begin{center}
    {\LARGE\bfseries\color{accent} Stanisław Szufa}\par
    \if\relax\detokenize{#2}\relax\else
      \vspace{3pt}{\large\color{muted} #2}\par
    \fi
    \ifheadercontact
    \vspace{5pt}
    {\small\color{muted}
    University of Geneva, Switzerland\par\vspace{1pt}
    \href{mailto:stanislaw.szufa@unige.ch}{stanislaw.szufa@unige.ch}\headersep
    +41 76 472 16 09%
    \if\relax\detokenize{#1}\relax\else\headersep#1\fi\par}
    \fi
  \end{center}
  \vspace{-6pt}
  {\color{accent!40}\rule{\linewidth}{0.5pt}}\par
  \vspace{2pt}}


\begin{document}

% ===============================================================
\docheader[\href{https://szufa.pl}{szufa.pl}\headersep
  \href{https://github.com/szufa}{github.com/szufa}]{Curriculum Vitae}

\textbf{Research area.} Computational social choice: algorithms and data for
collective decisions. I~develop methods to compare elections and to test voting
rules on real and synthetic data, I built the largest open dataset of
participatory budgeting votes, and I now study whether AI-written summaries
and LLM agents represent the opinions they aggregate.

% ===============================================================
\section{Academic Positions}
% ===============================================================
\entry{2026 -- now}{Postdoctoral Researcher}{University of Geneva}
  {Hosts: Roland Bouffanais and Didier Wernli. NLP pipelines for diplomatic
  and policy texts; computational models of decision-making in multilateral
  negotiations; LLM-based simulation of democratic deliberation.}
\entry{2024 -- 2026}{CNRS Postdoctoral Researcher}{LAMSADE, Université Paris Dauphine-PSL}
  {Host: Jérôme Lang. Social choice and machine learning; discovering hidden
structure in preference data.}
\entry{2021 -- 2024}{Research Assistant}{AGH University of Kraków}
  {ERC project PRAGMA led by Piotr Faliszewski. Maps of elections,
  participatory budgeting, development of \texttt{prefsampling}.}
\entry{2018 -- 2021}{Teaching and Research Assistant}{Jagiellonian University, Kraków}
  {Taught three undergraduate courses; founded Pabulib.}

% ===============================================================
\section{Education}
% ===============================================================
\entry{2017 -- 2024}{PhD in Computer Science}{AGH University of Kraków}
  {Doctoral studies at Jagiellonian University (2017--2023); thesis defended
  with distinction at AGH on 13 May 2024. Thesis: \emph{Map of Elections}.
  Advisor: Piotr Faliszewski.}
\entry{2015 -- 2017}{MSc in Computer Science}{Jagiellonian University}{}
\entry{2012 -- 2015}{BSc in Computer Science}{Jagiellonian University}{}
\vspace{4pt}

% ===============================================================
\section{Awards, Fellowships, and Grants}
% ===============================================================
\entry{2025}{Polish Prime Minister's Award}{}
  {For an outstanding doctoral dissertation (\emph{Map of Elections}).}
\entry{2023}{Pragnesh Jay Modi Best Student Paper Award}{AAMAS-2023}
  {For \emph{A Map of Diverse Synthetic Stable Roommates Instances}, with
  N.~Boehmer and K.~Heeger.}
\entry{2023}{START Fellowship}{Foundation for Polish Science}{}
\entry{2018}{Principal Investigator, NCN Preludium grant}{National Science Centre, Poland}
  {Project: \emph{Election Isomorphism Problem}. The project that started
  maps of elections.}
\vspace{4pt}

% ===============================================================
\section{Publications}
% ===============================================================
5 journal articles (\emph{Artificial Intelligence}, \emph{JAIR} $\times$2,
\emph{JCSS}, \emph{Theory and Decision}) and 28 full conference papers
(IJCAI $\times$11, AAMAS $\times$7, AAAI $\times$5, NeurIPS $\times$2,
ECAI $\times$2, Findings of EMNLP $\times$1).
The full list is provided as a separate document, with the same labels as below. Selected papers:
\vspace{3pt}
\paper{J4}{S. Szufa, N. Boehmer, R. Bredereck, P. Faliszewski,
  R. Niedermeier, P. Skowron, A. Slinko, N. Talmon. Drawing a Map of
  Elections. \emph{Artificial Intelligence} 343, 2025.}
\paper{C7}{S. Szufa, P. Faliszewski, Ł. Janeczko, M. Lackner, A. Slinko,
  K. Sornat, N. Talmon. How to Sample Approval Elections? \emph{IJCAI-2022}.}
\paper{C11}{P. Faliszewski, J. Flis, D. Peters, G. Pierczyński, P. Skowron,
  D. Stolicki, S. Szufa, N. Talmon. Participatory Budgeting: Data,
  Tools and Analysis. \emph{IJCAI-2023}.}
\paper{J1}{N. Boehmer, K. Heeger, S. Szufa. A Map of Diverse Synthetic
  Stable Matching Instances. \emph{JAIR} 79, 2024 (journal version of the AAMAS-2023
  Best Student Paper).}
\vspace{-2pt}
{\small\color{muted}Authors in computational social choice are usually listed
alphabetically; papers that deviate from alphabetical order list the lead
author first.}

% ===============================================================
\section{Software and Datasets}
% ===============================================================
\entry{2020 -- now}{Pabulib}{\href{https://pabulib.org}{pabulib.org}}
  {Founder. The largest open library of participatory
  budgeting data: 2{,}053 elections from 30 cities in 8 countries,
  8.6 million votes; cited in more than 100 papers.}
\entry{}{prefsampling}{Python library}
  {Co-developer, with Simon Rey. Sampling synthetic preference data under a wide range of
  statistical cultures.}
\entry{}{mapof}{Python library}
  {Developer. Computing distances between elections and drawing maps of
  elections.}

% ===============================================================
\section{Teaching and Supervision}
% ===============================================================
\entry{2018 -- 2021}{Undergraduate courses}{Jagiellonian University}
  {\emph{Introduction to Programming} (one year), \emph{Logic and Set Theory}
  and \emph{Linear Algebra} (one semester each); groups of about 20 students.}
\entry{2023 -- 2026}{Tutorials and workshops}{}
  {\emph{Experiments in Social Choice}, tutorial at IJCAI 2023, AAMAS 2023,
  and AAAI 2024.\\
  \emph{Computational Tools for Measuring Diversity, Agreement,
  and Polarization}, hands-on workshop at the Data Science Institute,
  University of Chicago, 2026.}
\entry{}{Co-supervision}{}
  {More than ten BSc and MSc students.}

% ===============================================================
\section{Service and Leadership}
% ===============================================================
\entry{2026 -- now}{Editorial board member}{Journal of Artificial Intelligence Research (JAIR)}{}
\entry{}{Program committee member}{NeurIPS, AAAI, AAMAS, IJCAI}{}
\entry{}{Journal reviewer}{JCSS, TCS, SCW}{}
\entry{2025 -- now}{Co-founder}{Nauka dla Demokracji (Science for Democracy)}{}
\entry{2021}{President}{Jagiellonian University PhD Students' Association}{}
\entry{2020}{Representative}{Polish Association of Doctoral Candidates}{}

% ===============================================================
\section{Research Visits}
% ===============================================================
\entry{2026}{University of Chicago}{Chicago, USA}{}
\entry{2025}{Northwestern University}{Evanston, USA}{}
\entry{2024}{University of Oxford}{Oxford, UK}{}
\entry{2023}{University of Warsaw}{Warsaw, Poland}{}
\entry{2023}{National University of Singapore}{Singapore}{}
\entry{2021}{TU Berlin}{Berlin, Germany}{}
% TODO (optional): add the host of each visit

% ===============================================================
\section{Invited Talks}
% ===============================================================
\entry{2025}{Dagstuhl Seminar 25401}{Societal Impact of Computational Social Choice, Germany}{}
\entry{2024}{Participatory Budgeting Workshop}{Society for Social Choice and Welfare}{}
\entry{2022}{Computational and Topological Social Choice Workshop}{Society for Social Choice and Welfare}{}

% ===============================================================
\section{Other Experience}
% ===============================================================
\entry{2020}{Optimization Consultant}{PCG Academia}
  {Integer linear programming models for real-world resource-allocation
  problems for a private-sector client.}

% ===============================================================
\section{Languages}
% ===============================================================
{\raggedright Polish (native), English (fluent), Spanish (intermediate),
French (conversational).\par}

\end{document}
